\section{总结与延伸}

本期以 Superpowers 项目为核心，系统讲解了 Vibe Coding 时代如何"驾驭"（Harness）AI Coding Agent，使其写出"又对又好又准"的代码。核心要点回顾：

\begin{enumerate}
    \item \textbf{Skills 加载机制}：通过 Prompt Injection 在 Codex 系统消息中暴露 Skill 目录，Agent 动态匹配并加载。
    \item \textbf{总控流程}：Using Superpowers $\rightarrow$ Brainstorming $\rightarrow$ Writing Plans $\rightarrow$ TDD Development $\rightarrow$ Verification。
    \item \textbf{TDD 纪律}：红-绿-重构三阶段，"亲眼看着测试失败"是不可跳过的第一步。
    \item \textbf{系统化调试}：四阶段闭环，三次未果则升级为架构问题。
    \item \textbf{Code Review 双向体系}：请求审查（Checklist 驱动）+ 接收审查（技术严谨性 + Push Back）。
    \item \textbf{DRY \& YAGNI}：贯穿始终的两大原则。
\end{enumerate}

Superpowers 的意义不仅仅是一组 Prompt，而是将一套经过验证的顶级软件工程实践"硬编码"为 AI 时代的规范操作系统。它代表了 \textbf{Harness Engineering}——通过 Skill 这根"缰绳"，驾驭越来越强的 Coding Agent，在效率与质量之间找到最优平衡。

\subsection{拓展阅读}

\begin{itemize}
    \item Superpowers GitHub 仓库
    \item Kent Beck, \textit{Test-Driven Development: By Example}
    \item Robert C. Martin, \textit{Clean Code}
    \item Martin Fowler, \textit{Refactoring: Improving the Design of Existing Code}
\end{itemize}

%% --- Fin del contenido --- %%

\end{document}
